\chapter*{Abstract}
\addcontentsline{toc}{chapter}{Abstract}
\vspace{0.2em}
{\singlespacing
\noindent
Lower-Layer Triggered Mobility (LTM), introduced in 3GPP Releases 18 and 19, replaces the slow RRC
handover with a single MAC Control Element (CE) to minimize mobility interruption for low-latency
applications. LTM achieves this speed-up by pre-caching candidate cell configurations, acquiring
candidate timing advances in advance, and executing cell switches without random access (RACH-less).
While prior work has noted that MAC CEs lack integrity protection, this report investigates the
architectural vulnerabilities introduced by these three speed-up mechanisms themselves. Two major
findings are identified through rigorous 3GPP specification analysis, both surviving prospective MAC CE
integrity protection because neither depends on message forgery.

\textbf{Finding 1} identifies that candidate timing advance acquisition relies on a physical-layer preamble
measurement bound only to a cell, preamble index, and time occasion---never to a UE identity. Furthermore,
no specification clause bounds how long the network may retain this measurement. A low-power RF relay
forwarding the preamble with differential delay can inflate the stored timing advance. Simulated against
open-source 5G stacks (srsRAN and OpenAirInterface), the measured decodability boundary spans 6--8
timing advance units (median 7); under flat-fading/AWGN channel conditions this reaches full, deterministic
uplink decode failure by 8 units, while realistic multipath saturates slightly more gradually across the
same range. This boundary is invariant across numerologies because the cyclic prefix scales identically
with the timing advance unit. The relay requires only a $+9.9$\,dB power margin under flat fading
($+5.3$\,dB under multipath) to succeed in 90\% of trials, placing it within an admissible radius equal
to or greater than the victim's distance.

\textbf{Finding 2} uncovers a protocol deadlock in the scheduling request (SR) procedure of TS~38.321
clause~5.4.4. Upon exhausting its retry counter (\texttt{sr-TransMax}), the UE unconditionally deallocates
its PUCCH, SRS, and all grants at depth \texttt{5>}. However, the recovery action (initiating random
access) sits at depth \texttt{6>}, gated by a condition that evaluates to false during an ongoing
RACH-less LTM switch. Across the 5400 legal 3GPP ASN.1 configurations, this deadlock is reachable in up to
77.4\% of configurations.

The two findings compose into a deterministic attack: when the UE applies the corrupted timing advance
during a RACH-less LTM switch, its misaligned transmit timing causes complete phase decorrelation on its
PUCCH Format~0 Scheduling Requests. Every SR attempt fails across all SNRs (a 23\,dB sweep). Counter
exhaustion occurs with near certainty under flat-fading/AWGN conditions ($P \ge 0.996$ at 7 units of
error) and remains high under realistic multipath ($P = 0.75$--$0.97$ depending on receive diversity and
the exact TA error), triggering unconditional resource teardown while recovery is blocked. We derive a
safe configuration rule ($\texttt{sr-TransMax} \times T_{\text{interval}}
\ge \texttt{t304}$) that closes the vulnerability, verify both findings against published vendor
configurations, and establish the Phase~II research roadmap.
}
